\chapter{Marco teórico} \label{chap:background}

En este capítulo se presentan los fundamentos teóricos que sustentan el desarrollo de esta tesis. Se abordan tres ejes principales: los sistemas de recomendación, el procesamiento de lenguaje natural con enfoque en el análisis de sentimientos basado en aspectos, y las técnicas de optimización metaheurística. Adicionalmente, se discuten las métricas de evaluación empleadas y los fundamentos de explicabilidad en sistemas de recomendación.

%% ============================================================
%% 2.1  SISTEMAS DE RECOMENDACIÓN
%% ============================================================
\section{Sistemas de recomendación} \label{sec:recsys}

% Definición formal del problema de recomendación.
% Matriz de utilidad R ∈ ℝ^{|U|×|I|}, donde la mayoría de entradas son desconocidas.
% Contextualizar con el dataset REST-MEX: |U|=9,582 usuarios, |I|=48 pueblos, sparsity=90.32%.

% --------------------------------------------------------------
\subsection{Filtrado colaborativo} \label{subsec:cf}

% Idea central: usuarios con comportamiento similar en el pasado tendrán preferencias similares en el futuro.

\subsubsection{Basado en Usuarios} \label{subsubsec:cf_memory}
% User-based CF: similaridad entre usuarios (coseno, Pearson).
% Item-based CF: similaridad entre ítems.
% Predicción como promedio ponderado por similaridad.
% Limitaciones: escalabilidad, sparsity, cold-start.

\subsubsection{Basado en Items} \label{subsubsec:cf_model}
% Factorización matricial: R ≈ P·Qᵀ, donde P ∈ ℝ^{|U|×k}, Q ∈ ℝ^{|I|×k}.
% SVD, SVD++, NMF.
% Función de pérdida con regularización:
%   min_{p,q} Σ_{(u,i)∈κ} (r_{ui} - p_u^T q_i)^2 + λ(||p_u||^2 + ||q_i||^2)
% Modelos baseline de la librería Surprise: KNNBasic, KNNWithMeans, SVD, SVD++, NMF, BaselineOnly.
% Ejemplo: en REST-MEX, con k factores latentes representando dimensiones como "cultura" o "gastronomía".

% --------------------------------------------------------------
\subsection{Filtrado basado en contenido} \label{subsec:cbf}

% Representación de ítems mediante vectores de características (perfil del ítem).
% Perfil del usuario como agregación de los perfiles de ítems consumidos.
% Similitud coseno entre perfil del usuario y candidatos.
% En este proyecto: características extraídas de reseñas textuales via ABSA (Sección \ref{sec:absa}).
% Ejemplo: perfil de aspectos de un pueblo = distribución de sentimientos por aspecto (servicio, gastronomía, naturaleza, ...).
% Ventaja sobre CF puro: mitiga cold-start de ítems.

% --------------------------------------------------------------
\subsection{Sistemas híbridos} \label{subsec:hybrid}

% Taxonomía de Burke (2002): weighted, switching, mixed, feature combination, cascade, feature augmentation, meta-level.
% Justificar la estrategia elegida para esta tesis (feature augmentation / weighted).
% Formalización: r̂_{ui} = α · r̂_{ui}^{CF} + (1-α) · r̂_{ui}^{CBF}, donde α es el parámetro de fusión.
% Extensión: α puede ser un vector de pesos optimizado por metaheurística (Sección \ref{sec:metaheuristic}).
% Ventaja: combina señales complementarias (patrones de rating + semántica textual).

% --------------------------------------------------------------
\subsection{Problema del arranque en frío y dispersión} \label{subsec:cold_start}

% Cold-start de usuario: nuevos usuarios sin historial → CBF puede explotar contenido textual.
% Cold-start de ítem: nuevos pueblos sin ratings → ABSA genera perfil desde reseñas.
% Sparsity: 90.32% en REST-MEX. Impacto en CF basado en memoria.
% Estrategias de mitigación: factorización matricial, incorporación de side information, sistemas híbridos.


%% ============================================================
%% 2.2  PROCESAMIENTO DE LENGUAJE NATURAL
%% ============================================================
\section{Procesamiento de lenguaje natural} \label{sec:nlp}

% Breve contexto: el texto de reseñas como fuente de información semántica complementaria a los ratings numéricos.

% --------------------------------------------------------------
\subsection{Representación de texto} \label{subsec:text_repr}

% Bag-of-Words, TF-IDF → limitaciones (pérdida de orden, alta dimensionalidad).
% Word embeddings (Word2Vec, GloVe) → representaciones densas.
% Embeddings contextuales (BERT, RoBERTa) → representaciones dependientes del contexto.

% --------------------------------------------------------------
\subsection{Arquitectura Transformer} \label{subsec:transformers}

% Mecanismo de self-attention: Attention(Q,K,V) = softmax(QKᵀ/√d_k)V.
% Multi-head attention, codificación posicional.
% Modelos pre-entrenados y fine-tuning para tareas downstream.
% Transfer learning: ventaja en dominios con datos etiquetados escasos (como turismo en español).


%% ============================================================
%% 2.3  ANÁLISIS DE SENTIMIENTOS BASADO EN ASPECTOS (ABSA)
%% ============================================================
\section{Análisis de sentimientos basado en aspectos} \label{sec:absa}

% Definición formal: dado un texto t, identificar un conjunto de tuplas {(a_i, s_i)} donde a_i es un aspecto y s_i ∈ {positivo, negativo, neutro}.
% Diferencia con análisis de sentimientos general: granularidad a nivel de aspecto, no de documento.
% Ejemplo REST-MEX: "La comida estaba deliciosa pero el hotel tenía mal servicio"
%   → {(gastronomía, positivo), (servicio, negativo)}.

% --------------------------------------------------------------
\subsection{Subtareas del ABSA} \label{subsec:absa_subtasks}

% Aspect Term Extraction (ATE): identificar menciones explícitas de aspectos en el texto.
% Aspect Category Detection (ACD): asignar categorías predefinidas (e.g., servicio, gastronomía, naturaleza).
% Aspect Sentiment Classification (ASC): determinar polaridad por aspecto.
% En esta tesis: ACD vía Zero-Shot + ASC vía modelo de sentimiento fine-tuned.

% --------------------------------------------------------------
\subsection{Clasificación Zero-Shot} \label{subsec:zero_shot}

% Paradigma: clasificar sin ejemplos etiquetados para las clases objetivo.
% Basado en Natural Language Inference (NLI): premisa=texto, hipótesis="Este texto habla sobre {aspecto}".
% Modelo utilizado: Recognai/zeroshot_selectra_medium (SELECTRA, entrenado en español).
% Ventaja: no requiere datos etiquetados específicos del dominio turístico para aspectos.
% Limitación: dependencia de la formulación de las hipótesis y del conjunto de labels candidatos.

% --------------------------------------------------------------
\subsection{Análisis de sentimientos con modelos pre-entrenados} \label{subsec:sentiment_models}

% Modelos para español: BETO, RoBERTuito, MarIA.
% Modelo utilizado: pysentimiento/robertuito-sentiment-analysis.
%   - Pre-entrenado en ~500M tweets en español.
%   - Fine-tuned para clasificación de sentimiento (POS, NEG, NEU).
% Pipeline: texto → segmentación en chunks → clasificación de aspecto (Zero-Shot) → clasificación de sentimiento → agregación.


%% ============================================================
%% 2.4  OPTIMIZACIÓN METAHEURÍSTICA
%% ============================================================
\section{Optimización metaheurística} \label{sec:metaheuristic}

% Problema: encontrar los pesos óptimos de fusión del sistema híbrido.
% Formulación: max_{θ} f(θ), donde θ = (α, hiperparámetros del modelo) y f es una métrica de ranking (e.g., NDCG@K).
% Espacio de búsqueda no convexo, no diferenciable → justificación de metaheurísticas.

% --------------------------------------------------------------
\subsection{Optimización por enjambre de partículas (PSO)} \label{subsec:pso}

% Algoritmo inspirado en comportamiento colectivo (bandadas, cardúmenes).
% Partícula i con posición x_i y velocidad v_i:
%   v_i(t+1) = ω·v_i(t) + c₁·r₁·(pbest_i - x_i(t)) + c₂·r₂·(gbest - x_i(t))
%   x_i(t+1) = x_i(t) + v_i(t+1)
% Parámetros: inercia ω, coeficientes cognitivo c₁ y social c₂.
% Balance exploración-explotación.

% --------------------------------------------------------------
\subsection{Optimización bayesiana} \label{subsec:bayesian_opt}

% Modelo sustituto (surrogate): Proceso Gaussiano GP(μ, k) que aproxima f(θ).
% Función de adquisición: Expected Improvement, UCB, Probability of Improvement.
%   EI(θ) = E[max(f(θ) - f(θ⁺), 0)]
% Ventaja: eficiente en evaluaciones costosas (cada evaluación requiere entrenar y evaluar el RecSys completo).
% Ciclo: evaluar → actualizar GP → optimizar adquisición → siguiente punto.

% --------------------------------------------------------------
\subsection{Evolución diferencial} \label{subsec:de}

% Operadores: mutación, cruce, selección.
% Mutación: v_i = x_{r1} + F·(x_{r2} - x_{r3}), donde F ∈ [0,2] es el factor de escala.
% Cruce binomial con probabilidad CR.
% Selección greedy: x_i(t+1) = u_i si f(u_i) ≥ f(x_i(t)), else x_i(t).
% Variantes: DE/rand/1, DE/best/1, DE/current-to-best/1.


%% ============================================================
%% 2.5  MÉTRICAS DE EVALUACIÓN
%% ============================================================
\section{Métricas de evaluación} \label{sec:metrics}

% Contexto: evaluación de sistemas Top-K; no basta RMSE/MAE del rating prediction.

% --------------------------------------------------------------
\subsection{Métricas de ranking} \label{subsec:ranking_metrics}

\subsubsection{Precision@K y Recall@K}
% Precision@K = |{ítems relevantes en top-K}| / K
% Recall@K = |{ítems relevantes en top-K}| / |{total ítems relevantes}|
% Trade-off precision-recall según K.

\subsubsection{NDCG@K}
% DCG@K = Σ_{i=1}^{K} (2^{rel_i} - 1) / log₂(i+1)
% NDCG@K = DCG@K / IDCG@K
% Penaliza ítems relevantes en posiciones bajas del ranking.
% Métrica principal en esta tesis.

\subsubsection{MRR (Mean Reciprocal Rank)}
% MRR = (1/|U|) Σ_{u=1}^{|U|} 1/rank_u
% Evalúa la posición del primer ítem relevante.

% --------------------------------------------------------------
\subsection{Métricas beyond-accuracy} \label{subsec:beyond_accuracy}

\subsubsection{Diversidad}
% Intra-list diversity: ILD = (2 / K(K-1)) Σ_{i<j} dist(item_i, item_j)
% Importancia en turismo: evitar recomendar K pueblos del mismo estado/tipo.

\subsubsection{Cobertura}
% Catalog coverage = |{ítems recomendados al menos una vez}| / |I|
% Relevante dado el desbalance geográfico del dataset (algunos pueblos con muchas más reseñas).


%% ============================================================
%% 2.6  EXPLICABILIDAD EN SISTEMAS DE RECOMENDACIÓN
%% ============================================================
\section{Explicabilidad en sistemas de recomendación} \label{sec:xai}

% Motivación: confianza del usuario, transparencia, regulación.
% Tipos de explicación: feature-based, example-based, textual.
% En esta tesis: explicabilidad basada en aspectos.
%   Ejemplo: "Se te recomienda Taxco porque valoras positivamente la cultura y la gastronomía,
%             aspectos altamente evaluados en este pueblo."
% Conexión ABSA → explicabilidad: los aspectos extraídos sirven como justificación interpretable.
% Post-hoc vs inherent interpretability.
